\section{Separation of births from earlier actual regions}
\label{sec:area_law_birth_separation}

A replacement born at a mark is compared with the final actual regions of
the older and peer identifiers of the same label, not with their births. The
comparison near the mark is reduced to a fixed normalized square, where a
compactness argument handles points away from the marks and an explicit
angular estimate handles points near a mark. The trim that follows the birth
then removes a square of radius comparable with $u^{\rho_0}$ about every mark
from every compared region.

Throughout, $\|\cdot\|$ is the maximum norm on $\mathbb R^2$. Normalize a
removal square to $Q=[-1,1]^2$, so that its center, corners and side
midpoints are the points of $\{-1,0,1\}^2$.

% Source: OpenAI2026AreaLaw, Section 11, 10-geometry.tex,
% geometry:angular-gap, lines 627--645.
\begin{theorem}[Angular estimate in the maximum norm]
    \label{thm:area_law_sup_angular_estimate}
    \lean{TNLean.PEPS.AreaLaw.Geometry.half_max_le_norm_smul_sub_smul,
        TNLean.PEPS.AreaLaw.Geometry.unitBox,
        TNLean.PEPS.AreaLaw.Geometry.one_le_norm_sub_of_disjoint_unitBox,
        TNLean.PEPS.AreaLaw.Geometry.norm_eq_one_of_mem_unitBox}
    \leanok
    Let $d,e\in\mathbb R^2$ satisfy $\|d\|=\|e\|=1$ and $\|d-e\|\ge1$, and let $r,s\ge0$. Then
    \begin{align}
        \|rd-se\| & \ge\tfrac12\max(r,s).
        \label{eq:area_law_sup_angular_estimate}
    \end{align}
    For $d=(d_1,d_2)\in\mathbb R^2$, set $B(d)=[\lfloor d_1\rfloor,\lceil d_1\rceil]\times[\lfloor d_2\rfloor,\lceil d_2\rceil]$.
    If $B(d)\cap B(e)=\emptyset$, then $\|d-e\|\ge1$. If $\|d\|=1$, then $B(d)\subseteq\{x:\|x\|=1\}$
    is the closed elementary half-side containing $d$ (or $\{d\}$ if $d\in\{-1,0,1\}^2\setminus\{0\}$).
\end{theorem}

\begin{proof}
    \leanok
    Assuming $s\le r$ and putting $D=\|rd-se\|$, the reverse triangle inequality gives
    $D\ge r-s$ and $r\le r\|d-e\|\le D+r-s$ gives $D\ge s$, so $D\ge\max(r-s,s)\ge r/2$.
    If $\|d-e\|<1$, then $\lfloor d_i\rfloor\le\lceil e_i\rceil$ and $\lfloor e_i\rfloor\le\lceil d_i\rceil$,
    so $(\max(\lfloor d_1\rfloor,\lfloor e_1\rfloor),\max(\lfloor d_2\rfloor,\lfloor e_2\rfloor))\in B(d)\cap B(e)$.
    For $\|d\|=1$, one coordinate is $\pm1$ (which $B(d)$ fixes) and the other lies in $[-1,1]$.
\end{proof}

% Source: OpenAI2026AreaLaw, Section 11, 10-geometry.tex,
% geometry:birth-separation, lines 573--659.
\begin{theorem}[Quarter sections of integer half-planes]
    \label{thm:area_law_quarter_sections}
    \lean{TNLean.PEPS.AreaLaw.Geometry.intHalfPlane,
        TNLean.PEPS.AreaLaw.Geometry.quarterRadialPoint,
        TNLean.PEPS.AreaLaw.Geometry.IsQuarterSectorial,
        TNLean.PEPS.AreaLaw.Geometry.isQuarterSectorial_intHalfPlane,
        TNLean.PEPS.AreaLaw.Geometry.IsQuarterSectorial.iUnion,
        TNLean.PEPS.AreaLaw.Geometry.IsQuarterSectorial.iInter}
    \leanok
    \uses{thm:area_law_sup_angular_estimate}
    A set $P\subseteq\mathbb R^2$ is \emph{quarter-sectorial at $q$} if:
    (i) for $0<\|x-q\|\le1/4$, $x\in P\iff q+\frac{x-q}{4\|x-q\|}\in P$; and
    (ii) for $\|d\|=1$, $q+d/4\in P\implies q+e/4\in P$ for all $e\in B(d)$.
    Quarter-sectoriality at $q$ is closed under arbitrary unions and intersections, and
    every closed half-plane $H_{a,b,n}=\{x\in\mathbb R^2:ax_1+bx_2\le n\}$ with
    $a,b,n\in\mathbb Z$ and $|a|+|b|\le2$ is quarter-sectorial at every $q\in\mathbb Z^2$.
\end{theorem}

\begin{proof}
    \leanok
    Set $f(x)=ax_1+bx_2-n$ and $k=f(q)\in\mathbb Z$. For $\|x-q\|\le1/4$,
    $|f(x)-k|\le(|a|+|b|)/4\le1/2$, so membership in $H_{a,b,n}$ is constant when $k\ne0$.
    When $k=0$, $f(q+w)=aw_1+bw_2$ is homogeneous, giving (i). On an open elementary half-side
    $d=(\pm1,t)$ with $t\in(m,m+1)$, $m\in\{-1,0\}$, $|a|+|b|\le2$ prevents $\pm a+bt$ from
    changing sign on $(m,m+1)$, so $\pm a+bt\le0$ extends to $t\in[m,m+1]$, giving (ii).
\end{proof}

% Source: OpenAI2026AreaLaw, Section 11, 10-geometry.tex,
% geometry:angular-gap, lines 627--645.
\begin{theorem}[Angular estimate at a contact mark]
    \label{thm:area_law_contact_angular_estimate}
    \lean{TNLean.PEPS.AreaLaw.Geometry.half_max_le_norm_sub_of_isQuarterSectorial}
    \leanok
    \uses{thm:area_law_sup_angular_estimate, thm:area_law_quarter_sections}
    If $P,P'\subseteq\mathbb R^2$ are quarter-sectorial at $q$ and
    $P\cap P'\cap\{z:\|z-q\|\le1/4\}\subseteq\{q\}$, then every $x\in P$ and $y\in P'$
    with $\|x-q\|,\|y-q\|\le1/4$ satisfy
    \begin{align}
        \|x-y\| & \ge\tfrac12\max(\|x-q\|,\|y-q\|).
        \label{eq:area_law_contact_angular_estimate}
    \end{align}
\end{theorem}

\begin{proof}
    \leanok
    If $x=q$ or $y=q$, (\ref{eq:area_law_contact_angular_estimate}) is immediate. Otherwise write
    $x=q+rd$, $y=q+se$ with $\|d\|=\|e\|=1$ and $r,s\in(0,1/4]$. By Theorem~\ref{thm:area_law_quarter_sections},
    $q+B(d)/4\subseteq P$ and $q+B(e)/4\subseteq P'$, so $P\cap P'\cap\{z:\|z-q\|=1/4\}=\emptyset$
    forces $B(d)\cap B(e)=\emptyset$. Theorem~\ref{thm:area_law_sup_angular_estimate} then gives
    $\|d-e\|\ge1$ and (\ref{eq:area_law_sup_angular_estimate}).
\end{proof}

% Source: OpenAI2026AreaLaw, Section 11, 10-geometry.tex, lines 393--417
% and geometry:birth-separation, lines 606--625.
\begin{definition}[Normalized replacement fan and exterior trapezoids]
    \label{def:area_law_normalized_repair_fan}
    \lean{TNLean.PEPS.AreaLaw.Geometry.repairSlotFacets,
        TNLean.PEPS.AreaLaw.Geometry.repairNormTriangle,
        TNLean.PEPS.AreaLaw.Geometry.repairExteriorTrapezoid,
        TNLean.PEPS.AreaLaw.Geometry.RepairSlotConsecutive}
    \leanok
    \uses{def:area_law_repair_square}
    For each of the $8$ elementary wedges $s$ of $Q=[-1,1]^2$, let $o_s$ be the outward unit
    normal of its side and $T_s$ the closed triangle joining $0$ to its side segment.
    For $L\in\{1,2\}$, $LT_s$ is defined by $o_s\cdot x\le L$ and the two homogeneous radial-edge
    inequalities; the exterior trapezoid is $E_s=\{x\in2T_s:o_s\cdot x\ge1\}=2T_s\setminus\operatorname{int}Q$.
    Two wedges are consecutive when they share a radial edge.
\end{definition}

% Source: OpenAI2026AreaLaw, Section 11, 10-geometry.tex,
% geometry:repair-induction, lines 497--504, and
% geometry:birth-separation, lines 616--625.
\begin{theorem}[Contacts in the normalized replacement fan]
    \label{thm:area_law_normalized_fan_contacts}
    \lean{TNLean.PEPS.AreaLaw.Geometry.isQuarterSectorial_repairNormTriangle,
        TNLean.PEPS.AreaLaw.Geometry.isQuarterSectorial_repairExteriorTrapezoid,
        TNLean.PEPS.AreaLaw.Geometry.isCompact_repairNormTriangle,
        TNLean.PEPS.AreaLaw.Geometry.isCompact_repairExteriorTrapezoid,
        TNLean.PEPS.AreaLaw.Geometry.repairNormTriangle_inter_exteriorTrapezoid,
        TNLean.PEPS.AreaLaw.Geometry.repairNormTriangle_inter_eq_zero,
        TNLean.PEPS.AreaLaw.Geometry.cellFanPolygon_region_subset_repairNormTriangle}
    \leanok
    \uses{def:area_law_normalized_repair_fan, thm:area_law_quarter_sections}
    The sets $LT_s$ ($L\in\{1,2\}$) and $E_s$ (Definition~\ref{def:area_law_normalized_repair_fan})
    are compact and quarter-sectorial at every $q\in\mathbb Z^2$. For $s\ne t$,
    $T_s\cap E_t\subseteq\mathbb Z^2$, and if $s,t$ are non-consecutive then $T_s\cap T_t=\{0\}$.
    Every actual fan triangle of $Q$ lies in $T_s$.
\end{theorem}

\begin{proof}
    \leanok
    Each facet has $a,b,n\in\mathbb Z$ with $|a|+|b|\le2$, so Theorem~\ref{thm:area_law_quarter_sections}
    gives quarter-sectoriality, and boundedness gives compactness. The intersections $T_s\cap E_t$
    and $T_s\cap T_t$ reduce to finite linear systems over pairs $(s,t)$, and each actual fan triangle
    is the convex hull of three vertices in $T_s$.
\end{proof}

% Source: OpenAI2026AreaLaw, Section 11, 10-geometry.tex,
% geometry:birth-separation, lines 616--645.
\begin{theorem}[Compactness away from the marks]
    \label{thm:area_law_contact_dichotomy}
    \lean{TNLean.PEPS.AreaLaw.Geometry.exists_pos_le_norm_sub_away_from_intPoints,
        TNLean.PEPS.AreaLaw.Geometry.exists_pos_contact_dichotomy,
        TNLean.PEPS.AreaLaw.Geometry.exists_pos_forall_of_finite}
    \leanok
    \uses{thm:area_law_contact_angular_estimate}
    If $X,Y\subseteq\mathbb R^2$ are compact, quarter-sectorial at every $q\in\mathbb Z^2$,
    and satisfy $X\cap Y\subseteq\mathbb Z^2$, then there exists $\eta>0$ such that for all
    $x\in X$ and $y\in Y$, either $\|x-y\|\ge\eta$, or there exists $q\in\mathbb Z^2$ with
    $\|x-q\|,\|y-q\|<1/4$ and $\|x-y\|\ge\frac12\|y-q\|$.
\end{theorem}

\begin{proof}
    \leanok
    The set of $(x,y)\in X\times Y$ not lying together in any open ball $\{z:\|z-q\|<1/4\}$
    ($q\in\mathbb Z^2$) is compact and disjoint from the diagonal (since $X\cap Y\subseteq\mathbb Z^2$),
    so $\|x-y\|$ attains a minimum $\eta>0$ on it; on each ball $\{z:\|z-q\|<1/4\}$,
    $X\cap Y$ meets the ball at most at $q$ and Theorem~\ref{thm:area_law_contact_angular_estimate}
    gives (\ref{eq:area_law_contact_angular_estimate}).
\end{proof}

% Source: OpenAI2026AreaLaw, Section 11, 10-geometry.tex,
% geometry:birth-separation, lines 573--659.
\begin{theorem}[Birth separation of a replacement from local comparison
        regions]
    \label{thm:area_law_replacement_birth_separation_local}
    \lean{TNLean.PEPS.AreaLaw.Geometry.repairNormRunRegion,
        TNLean.PEPS.AreaLaw.Geometry.repairSameLabelBound,
        TNLean.PEPS.AreaLaw.Geometry.exists_repairSameLabel_dichotomy,
        TNLean.PEPS.AreaLaw.Geometry.exists_repairBirth_separation}
    \leanok
    \uses{thm:area_law_normalized_fan_contacts,
        thm:area_law_contact_dichotomy, thm:area_law_repair_margins,
        def:area_law_repair_children}
    There exists $c_1>0$ such that for any exterior wedge labels, any maximal cyclic run
    $C$ of equal label with $B_C=\bigcup_{s\in C}T_s$, and
    \begin{align}
        W_C & = \Bigl(\bigcup_{\lambda(t)\ne\lambda(C)}E_t\Bigr)\cup\Bigl(\bigcup_{C'\ne C,\,\lambda(C')=\lambda(C)}B_{C'}\Bigr),
        \label{eq:area_law_local_comparison_union}
    \end{align}
    if $v\in\mathbb R^2$, $u\ge1$, and $Y\subseteq\mathbb R^2$ satisfies
    (1) $(y-v)/u\in W_C$ whenever $y\in Y$ and $\|y-v\|<2u$, and
    (2) $Y\cap\bigcup_{p\in\{-1,0,1\}^2}Q(v+up,u)=\emptyset$,
    then $\|x-y\|\ge c_1u^{\rho_0}$ for all $x\in v+uB_C$ and $y\in Y$.
\end{theorem}

\begin{proof}
    \leanok
    By Theorems~\ref{thm:area_law_normalized_fan_contacts} and \ref{thm:area_law_contact_dichotomy}
    applied to each pair $(T_s,E_t)$ or $(T_s,T_{t'})$ in $B_C\times W_C$ (where $t'\notin C$ is
    non-consecutive to $s\in C$), there is a uniform $\eta\in(0,1]$; set $c_1=\min(\eta,1/4)$.
    For $x\in v+uB_C$ and $y\in Y$: if $\|y-v\|\ge2u$, then $\|x-y\|\ge u\ge u^{\rho_0}$;
    if $\|y-v\|<2u$, rescaling by $u$ gives either $\|x-y\|\ge\eta u\ge c_1u^{\rho_0}$ or some
    $p\in\{-1,0,1\}^2$ with $q=v+up$ and $\|x-y\|\ge\frac12\|y-q\|>\frac12 d(u)\ge\frac14 u^{\rho_0}$
    by hypothesis (2) and Theorem~\ref{thm:area_law_repair_margins}.
\end{proof}

% Source: OpenAI2026AreaLaw, Section 11, 10-geometry.tex,
% geometry:initial-stars, lines 332--349; geometry:repair-induction,
% lines 486--494; geometry:birth-separation, lines 584--615.
\begin{theorem}[Same-label regions in an isolated star]
    \label{thm:area_law_star_same_label_bound}
    \lean{TNLean.PEPS.AreaLaw.Geometry.interior_repairNormTriangle_inter_eq_empty,
        TNLean.PEPS.AreaLaw.Geometry.interior_repairCurrentFanTriangle_inter_eq_empty,
        TNLean.PEPS.AreaLaw.Geometry.IsIsolatedStar.exists_repairCurrentSectorAssignment_eq,
        TNLean.PEPS.AreaLaw.Geometry.IsIsolatedStar.exists_mem_repairExteriorTrapezoid}
    \leanok
    \uses{thm:area_law_normalized_fan_contacts}
    If $v+[-\rho,\rho]^2$ is an isolated star of a labelled region family $(R_i)$ and
    $0<6u\le\rho$, then the working-wedge assignment at half-side $6u$ equals the star's
    wedge assignment, and every $w\in R_i$ with $u<\|w-v\|<2u$ satisfies $(w-v)/u\in E_t$
    for some wedge $t$ assigned to $i$.
\end{theorem}

\begin{proof}
    \leanok
    Since $\operatorname{int}T_t\cap\operatorname{int}T_s=\emptyset$ for $s\ne t$, a point in the
    open working triangle of $t$ lies in the chosen region for $t$ and hence in the star triangle
    of $t$. Near $v$, $\overline{R_i}$ is the union of the star triangles of wedges $t$ assigned to $i$;
    restricting to $u<\|w-v\|<2u$ and rescaling by $u$ places $(w-v)/u$ in
    $\{x\in2T_t:\|x\|=o_t\cdot x>1\}\subseteq E_t$.
\end{proof}

% Source: OpenAI2026AreaLaw, Section 11, 10-geometry.tex,
% geometry:birth-separation, lines 573--659, from
% geometry:repair-induction, lines 471--485.
\begin{theorem}[Birth separation of replacements, given the uniform repair
        induction]
    \label{thm:area_law_replacement_birth_separation_conditional}
    \lean{TNLean.PEPS.AreaLaw.Geometry.repairRegionLift,
        TNLean.PEPS.AreaLaw.Geometry.repairFinalRegion,
        TNLean.PEPS.AreaLaw.Geometry.IsRepairBirthSeparated,
        TNLean.PEPS.AreaLaw.Geometry.exists_isRepairBirthSeparated_of_isUniformRepairInduction,
        TNLean.PEPS.AreaLaw.Geometry.exists_isRepairBirthSeparated_of_uniformRepairInductionStatement}
    \leanok
    \uses{thm:area_law_replacement_birth_separation_local,
        thm:area_law_star_same_label_bound,
        def:area_law_uniform_repair_conclusion,
        lem:area_law_repair_metric_choice}
    Assuming the uniform repair induction (Definition~\ref{def:area_law_uniform_repair_conclusion}),
    there exist $c_1>0$ and $T$ such that for every $S_*\ge T$, $C_0\ge2$, and sufficiently large $k_0$,
    every replacement born at a parent mark of scale $S$ has closed birth at distance
    $\ge c_1d(S)^{\rho_0}$ from the final actual region $R_y^{(\infty,j+1)}$ of every other
    identifier $y$ of generation $j+1$ with the same label.
\end{theorem}

\begin{proof}
    \leanok
    Apply Theorem~\ref{thm:area_law_replacement_birth_separation_local} with $u=d(S)\ge S_*\ge1$
    to $Y=R_y^{(\infty,j+1)}$. Within $2u$ of the parent mark $v$, an older $y$ avoids $Q(v,S)$
    and lies in $uE_t$ for wedges $t$ of its label (opposite to the run's exterior label) by
    Theorem~\ref{thm:area_law_star_same_label_bound}; a peer $y$ at $v$ lies in $uB_{C'}$ for
    another run $C'$ of the same exterior label; and a peer at another parent $v'$ lies in
    $Q(v',S')$ at distance $\ge2u$ from $v$ by mark separation and Lemma~\ref{lem:area_law_repair_metric_choice}.
    Generation $j+1$ removes $\bigcup_{p\in\{-1,0,1\}^2}Q(v+up,u)$ from $Y$.
\end{proof}

For an initial belt tile in a belt cell of side $u=t_k$, normalize the cell to $[0,4]^2$
(in units of $u/4$). Cells of the same and adjacent layers have sides in $\{2,4,8\}$,
corners in $(2\mathbb Z)^2$, and sides and fan edges on the arrangement
\begin{align}
    x_1=m,\qquad x_2=m,\qquad x_1+x_2=2m,\qquad x_1-x_2=2m,
    \qquad m\in\mathbb Z,
    \label{eq:area_law_belt_arrangement}
\end{align}
whose faces are the \emph{atoms} $A_{c,s}=(2c_1+1,2c_2+1)+T_s$ ($c\in\mathbb Z^2$,
Definition~\ref{def:area_law_normalized_repair_fan}).

% Source: OpenAI2026AreaLaw, Section 11, 10-geometry.tex,
% geometry:initial-stars, lines 352--359, and geometry:birth-separation,
% lines 597--625.
\begin{theorem}[Atoms of the normalized belt arrangement]
    \label{thm:area_law_belt_atoms}
    \lean{TNLean.PEPS.AreaLaw.Geometry.IsBeltArrangementPoint,
        TNLean.PEPS.AreaLaw.Geometry.IsBeltFacet,
        TNLean.PEPS.AreaLaw.Geometry.exists_int_eq_of_isBeltFacet,
        TNLean.PEPS.AreaLaw.Geometry.beltAtomCenter,
        TNLean.PEPS.AreaLaw.Geometry.beltAtom,
        TNLean.PEPS.AreaLaw.Geometry.isCompact_beltAtom,
        TNLean.PEPS.AreaLaw.Geometry.isQuarterSectorial_beltAtom,
        TNLean.PEPS.AreaLaw.Geometry.not_isBeltArrangementPoint_of_mem_interior,
        TNLean.PEPS.AreaLaw.Geometry.exists_beltAtom_ray,
        TNLean.PEPS.AreaLaw.Geometry.exists_beltAtom_subset,
        TNLean.PEPS.AreaLaw.Geometry.exists_segment_subset_beltAtom_inter,
        TNLean.PEPS.AreaLaw.Geometry.beltAtom_axis_segments}
    \leanok
    \uses{def:area_law_normalized_repair_fan, thm:area_law_quarter_sections}
    Every atom $A_{c,s}$ is compact, convex, and quarter-sectorial at every $q\in\mathbb Z^2$,
    and $\operatorname{int}A_{c,s}$ misses all lines of (\ref{eq:area_law_belt_arrangement}). Moreover:
    \begin{enumerate}
        \item if $P\subseteq\mathbb R^2$ is closed and convex with $\operatorname{int}P\ne\emptyset$
              and $\partial P$ contained in the lines (\ref{eq:area_law_belt_arrangement}), then every
              $x\in P$ lies in an atom $A_{c,s}\subseteq P$;
        \item two atoms meeting outside $\mathbb Z^2$ share a nondegenerate segment;
        \item if $q\in\mathbb Z^2$ has $q_1+q_2$ odd, every atom containing $q$ contains
              $[q,q+\sigma e_1]$ and $[q,q+\tau e_2]$ for some $\sigma,\tau\in\{\pm1\}$.
    \end{enumerate}
\end{theorem}

\begin{proof}
    \leanok
    Each atom is cut out by three half-planes of (\ref{eq:area_law_belt_arrangement}) with
    $|a|+|b|\le2$, so Theorem~\ref{thm:area_law_quarter_sections} applies, nonparallel lines meet
    in $\mathbb Z^2$, and $\operatorname{int}A_{c,s}$ avoids (\ref{eq:area_law_belt_arrangement}).
    For (1), a ray from $x\in P$ to $c\in\operatorname{int}P$ enters an atom whose connected
    interior meets $\operatorname{int}P$ and misses $\partial P$, hence lies in $P$.
    For (2), facets tight at a common point outside $\mathbb Z^2$ lie on a single line of
    (\ref{eq:area_law_belt_arrangement}), giving a shared segment. For (3), $q$ with $q_1+q_2$ odd
    is a side midpoint of $2\mathbb Z^2$, where each incident atom contains both axis unit segments.
\end{proof}

% Source: OpenAI2026AreaLaw, Section 11, 10-geometry.tex,
% geometry:birth-separation, lines 597--645.
\begin{theorem}[Uniform dichotomy for atoms]
    \label{thm:area_law_belt_atom_dichotomy}
    \lean{TNLean.PEPS.AreaLaw.Geometry.exists_pos_common_contact_dichotomy,
        TNLean.PEPS.AreaLaw.Geometry.exists_beltAtom_dichotomy}
    \leanok
    \uses{thm:area_law_belt_atoms}
    There exists $\eta\in(0,1]$ such that for any atoms $A,B$ with $A\cap B\subseteq\mathbb Z^2$
    and any $x\in A$, $y\in B$, either $\|x-y\|\ge\eta$, or there exists $q\in A\cap B$ with
    $\|x-q\|,\|y-q\|<1/4$ and $\|x-y\|\ge\frac12\|y-q\|$.
\end{theorem}

\begin{proof}
    \leanok
    \uses{thm:area_law_contact_angular_estimate}
    By Theorems~\ref{thm:area_law_belt_atoms} and \ref{thm:area_law_contact_angular_estimate}, each pair
    $(A,B)$ with $A\cap B\subseteq\mathbb Z^2$ satisfies the dichotomy for some $\eta_{A,B}>0$;
    modulo $(2\mathbb Z)^2$-translations preserving (\ref{eq:area_law_belt_arrangement}), only finitely
    many pairs satisfy $\operatorname{dist}(A,B)<1$.
\end{proof}

% Source: OpenAI2026AreaLaw, Section 11, 10-geometry.tex,
% geometry:birth-separation, lines 592--604, with the scale comparison,
% lines 193--207.
\begin{theorem}[Initial regions near a belt cell]
    \label{thm:area_law_belt_cell_local_atoms}
    \lean{TNLean.PEPS.AreaLaw.Geometry.beltCellOrigin,
        TNLean.PEPS.AreaLaw.Geometry.beltNormalize,
        TNLean.PEPS.AreaLaw.Geometry.exists_beltAtom_subset_cellFanPolygon,
        TNLean.PEPS.AreaLaw.Geometry.exists_beltAtom_subset_closure_dyadicCell,
        TNLean.PEPS.AreaLaw.Geometry.exists_beltAtom_subset_initialBirthRegion}
    \leanok
    \uses{thm:area_law_belt_atoms, def:area_law_initial_birth_regions}
    Under normalization of a layer-$k$ belt cell to $[0,4]^2$, every fan triangle and closed
    fine cell of any layer $h$ with $\ell_h\ge\ell_k-1$ satisfies Theorem~\ref{thm:area_law_belt_atoms}(1).
    Consequently, for $C_0\ge2$ and $k\ge k_0\ge50{,}000{,}000$, if $p$ lies in the closed belt
    cell and $w$ lies in an initial birth region $X_y^0$ (Definition~\ref{def:area_law_initial_birth_regions})
    with $\|p-w\|<t_k$, then the normalized point $\tilde w$ lies in an atom contained in the
    normalized birth region $\tilde X_y^0$.
\end{theorem}

\begin{proof}
    \leanok
    \uses{thm:area_law_local_layer_indices, thm:area_law_adjacent_fine_meshes,
        thm:area_law_primary_fine_cell_cover,
        thm:area_law_fan_frontier_primitives, thm:area_law_cell_fan_cover,
        thm:area_law_cell_fan_regularity}
    By Theorems~\ref{thm:area_law_adjacent_fine_meshes} and \ref{thm:area_law_fan_frontier_primitives},
    cells and fan triangles at scale index $\ell_h\ge\ell_k-1$ have frontiers on
    (\ref{eq:area_law_belt_arrangement}). By Theorem~\ref{thm:area_law_local_layer_indices}, any $w$
    within $t_k$ of the cell belongs to layer $h\in\{k-1,k,k+1\}$ (so $\ell_h\ge\ell_k-1$), and
    $X_y^0$ is a union of closed fine cells (Theorem~\ref{thm:area_law_primary_fine_cell_cover}) or
    fan triangles (Theorems~\ref{thm:area_law_cell_fan_cover} and \ref{thm:area_law_cell_fan_regularity}).
\end{proof}

% Source: OpenAI2026AreaLaw, Section 11, 10-geometry.tex,
% geometry:birth-separation, lines 606--615, with the belt construction,
% lines 299--330.
\begin{theorem}[Contacts of an initial belt tile]
    \label{thm:area_law_initial_belt_contacts}
    \lean{TNLean.PEPS.AreaLaw.Geometry.initialBeltRun,
        TNLean.PEPS.AreaLaw.Geometry.exists_eq_initialBeltRun_of_inter_interior,
        TNLean.PEPS.AreaLaw.Geometry.not_segment_subset_initialBirthRegion_inter,
        TNLean.PEPS.AreaLaw.Geometry.add_mem_segment_of_odd,
        TNLean.PEPS.AreaLaw.Geometry.mem_beltCellMarks_of_beltNormalize_eq,
        TNLean.PEPS.AreaLaw.Geometry.beltNormalize_symm_mem_beltCellMarks}
    \leanok
    \uses{thm:area_law_belt_cell_local_atoms, def:area_law_belt_cell_marks}
    For $C_0\ge2$ and $k\ge k_0\ge50{,}000{,}000$, let $x$ be an initial belt tile in a layer-$k$
    belt cell and $y\ne x$ an initial identifier with $\lambda(y)=\lambda(x)$. If $q\in\mathbb Z^2$
    (in cell-normalized coordinates) lies in $\tilde X_x^0$ and in an atom $B\subseteq\tilde X_y^0$,
    then $q\in\{0,2,4\}^2$ is the normalized position of one of the $9$ marks of the cell
    (Definition~\ref{def:area_law_belt_cell_marks}).
\end{theorem}

\begin{proof}
    \leanok
    \uses{thm:area_law_initial_interface_colors, thm:area_law_initial_open_disjoint,
        thm:area_law_initial_birth_regularity, thm:area_law_fan_run_geometric_contact,
        thm:area_law_square_frontier_side, thm:area_law_fan_own_side_contact}
    By Theorem~\ref{thm:area_law_initial_interface_colors}, $\tilde X_x^0$ and $\tilde X_y^0$ share no
    nondegenerate segment. If $\operatorname{int}X_y^0$ meets the open belt cell,
    Theorems~\ref{thm:area_law_initial_open_disjoint} and \ref{thm:area_law_fan_run_geometric_contact}
    make $y$ another same-label run of the cell meeting $\tilde X_x^0$ only at the centre $q=(2,2)$.
    Otherwise $\tilde X_y^0$ avoids $(0,4)^2$, so $q$ lies on $\partial[0,4]^2$ along a fan-triangle side
    of $\tilde X_x^0$ with endpoints in $\{0,2,4\}^2$ (Theorems~\ref{thm:area_law_square_frontier_side}
    and \ref{thm:area_law_fan_own_side_contact}); if $q_1+q_2$ were odd,
    Theorem~\ref{thm:area_law_belt_atoms}(3) would give a shared unit segment in $\tilde X_x^0\cap B$,
    contradicting segment-disjointness.
\end{proof}

% Source: OpenAI2026AreaLaw, Section 11, 10-geometry.tex,
% geometry:birth-separation, lines 573--659.
\begin{theorem}[Birth separation of an initial belt tile from local
        comparison regions]
    \label{thm:area_law_initial_belt_birth_separation_local}
    \lean{TNLean.PEPS.AreaLaw.Geometry.exists_initialBeltRun_birth_separation}
    \leanok
    \uses{def:area_law_initial_birth_regions, def:area_law_initial_open_regions,
        def:area_law_belt_cell_marks}
    There exists $c_1>0$ such that for $C_0\ge2$, $k\ge k_0\ge50{,}000{,}000$, an initial belt tile
    $x$ in a belt cell of side $u=t_k$, any primary or initial belt tile $y\ne x$ with
    $\lambda(y)=\lambda(x)$, and any $Y\subseteq\operatorname{int}X_y^0$ at distance $>u^{\rho_0}/4$
    from each of the $9$ marks of the cell (Definition~\ref{def:area_law_belt_cell_marks}),
    every $p\in X_x^0$ and $w\in Y$ satisfy $\|p-w\|\ge c_1u^{\rho_0}$.
\end{theorem}

\begin{proof}
    \leanok
    \uses{thm:area_law_belt_atom_dichotomy, thm:area_law_belt_cell_local_atoms,
        thm:area_law_initial_belt_contacts, thm:area_law_belt_atoms}
    Set $c_1=\min(\eta/4,1/8)$ with $\eta$ from Theorem~\ref{thm:area_law_belt_atom_dichotomy}.
    If $\|p-w\|\ge u$, then $\|p-w\|\ge u^{\rho_0}\ge c_1u^{\rho_0}$. If $\|p-w\|<u$,
    Theorem~\ref{thm:area_law_belt_cell_local_atoms} places the normalized points $\tilde p,\tilde w$
    in atoms $A\subseteq\tilde X_x^0$ and $B\subseteq\tilde X_y^0$ with $A\cap B\subseteq\mathbb Z^2$
    (Theorem~\ref{thm:area_law_belt_atoms}(2)). By Theorem~\ref{thm:area_law_belt_atom_dichotomy},
    either $\|p-w\|\ge\frac{\eta}{4}u\ge c_1u^{\rho_0}$, or there is a mark $v$ (by
    Theorem~\ref{thm:area_law_initial_belt_contacts}) with $\|p-w\|\ge\frac12\|w-v\|>\frac18 u^{\rho_0}$.
\end{proof}

% Source: OpenAI2026AreaLaw, Section 11, 10-geometry.tex,
% geometry:birth-separation, lines 573--659, with the order of
% identifiers, lines 563--571.
\begin{theorem}[Birth separation of the initial belt tiles]
    \label{thm:area_law_initial_belt_birth_separation}
    \lean{TNLean.PEPS.AreaLaw.Geometry.InitialRegionIndex.beltScaleIndex,
        TNLean.PEPS.AreaLaw.Geometry.IsInitialBirthSeparated,
        TNLean.PEPS.AreaLaw.Geometry.exists_isInitialBirthSeparated}
    \leanok
    \uses{def:area_law_repair_stage}
    There exists $c_1>0$ such that for every finite domain, cut, origin, $C_0\ge2$,
    $k_0\ge50{,}000{,}000$, and stopping scale $S_*$, the closed birth region $\mathcal B=X_x^0$
    of any initial belt tile $x$ of side $u=t_k$ has distance $\ge c_1u^{\rho_0}$ from the
    final actual region $R_y^{(\infty,0)}$ of every primary or other initial belt tile $y$ with
    $\lambda(y)=\lambda(x)$.
\end{theorem}

\begin{proof}
    \leanok
    \uses{thm:area_law_initial_belt_birth_separation_local,
        thm:area_law_uniform_separated_initial_marks, thm:area_law_repair_margins,
        thm:area_law_local_layer_indices, thm:area_law_adjacent_fine_meshes}
    Every $R_y^{(\infty,0)}\subseteq\operatorname{int}X_y^0$ avoids the generation-$0$ removal squares
    around the $9$ marks of the belt cell. Each such mark has initial scale $S\ge u/2$
    (Theorems~\ref{thm:area_law_uniform_separated_initial_marks}, \ref{thm:area_law_local_layer_indices},
    and \ref{thm:area_law_adjacent_fine_meshes}) and removal half-side $d(S)\ge S^{\rho_0}/2\ge u^{\rho_0}/4$
    (Theorem~\ref{thm:area_law_repair_margins}), so Theorem~\ref{thm:area_law_initial_belt_birth_separation_local}
    applies to $Y=R_y^{(\infty,0)}$.
\end{proof}

% Source: OpenAI2026AreaLaw, Section 11, 10-geometry.tex,
% geometry:birth-separation, lines 573--659, from
% geometry:repair-induction, lines 471--485.
\begin{theorem}[Birth-template separation, given the uniform repair
        induction]
    \label{thm:area_law_birth_separation_conditional}
    \lean{TNLean.PEPS.AreaLaw.Geometry.IsInitialBirthSeparated.mono,
        TNLean.PEPS.AreaLaw.Geometry.IsRepairBirthSeparated.mono,
        TNLean.PEPS.AreaLaw.Geometry.exists_birthSeparation_of_uniformRepairInductionStatement,
        TNLean.PEPS.AreaLaw.Geometry.exists_birthSeparation_of_initialIsolatedStars}
    \leanok
    \uses{def:area_law_uniform_repair_conclusion}
    Assuming the uniform repair induction (Definition~\ref{def:area_law_uniform_repair_conclusion})
    or the initial isolated-star lemma, there exist $c_1>0$ and $T$ such that for every
    $S_*\ge T$, $C_0\ge2$, and sufficiently large $k_0$, both
    Theorems~\ref{thm:area_law_replacement_birth_separation_conditional} and
    \ref{thm:area_law_initial_belt_birth_separation} hold with constant $c_1$.
\end{theorem}

\begin{proof}
    \leanok
    \uses{thm:area_law_replacement_birth_separation_conditional,
        thm:area_law_initial_belt_birth_separation}
    Take $c_1$ as the minimum of the constants from
    Theorems~\ref{thm:area_law_replacement_birth_separation_conditional} and
    \ref{thm:area_law_initial_belt_birth_separation}.
\end{proof}

Combining Theorem~\ref{thm:area_law_initial_stars} with
Theorem~\ref{thm:area_law_birth_separation_conditional} yields
Theorem~\ref{thm:area_law_birth_separation}~\cite{gap:arealaw2d_birth_separation_scope}.

% Source: OpenAI2026AreaLaw, Section 11, 10-geometry.tex,
% geometry:birth-separation, lines 573--582.
\begin{theorem}[Birth-template separation]
    \label{thm:area_law_birth_separation}
    \lean{TNLean.PEPS.AreaLaw.Geometry.exists_birthSeparation}
    \leanok
    \uses{def:area_law_repair_stage}
    There exists $c_1>0$ such that for every $C_0\ge2$ and sufficiently large $S_*$ and $k_0$,
    the closed birth region $\mathcal B$ of any initial belt tile (of scale $u=t_k$) or replacement
    tile (of removal half-side $u$) has distance $\ge c_1u^{\rho_0}$ from the final actual region
    of every earlier or peer identifier of the same label.
\end{theorem}

\begin{proof}
    \leanok
    \uses{thm:area_law_birth_separation_conditional,
        thm:area_law_initial_stars}
    Immediate from Theorems~\ref{thm:area_law_initial_stars} and
    \ref{thm:area_law_birth_separation_conditional}.
\end{proof}

% Auxiliary to OpenAI2026AreaLaw, Section 11, 10-geometry.tex,
% geometry:birth-separation, lines 565--591, and its use, lines 784--789.
% Revision: adc7f1241b42e322a6451854ab7e4b4c146bf78a.
\begin{theorem}[Birth separation at a common generation]
    \label{thm:area_law_birth_separation_common_stage}
    \lean{TNLean.PEPS.AreaLaw.Geometry.repairStage_birth_separation_of_rank_le}
    \leanok
    \uses{def:area_law_repair_stage,def:area_law_repair_inclusions,
        def:area_law_repair_birth_rank,def:area_law_repair_characteristic_length,
        thm:area_law_initial_belt_birth_separation,
        thm:area_law_replacement_birth_separation_conditional}
    In the repair family $\mathcal D_j$ (Definition~\ref{def:area_law_repair_stage}), fix
    $c_1\in\mathbb R$ and assume the separation bounds of
    Theorems~\ref{thm:area_law_initial_belt_birth_separation} and
    \ref{thm:area_law_replacement_birth_separation_conditional} with constant $c_1$.
    For any $N\ge0$ and $x,y\in I_N$ with $\operatorname{rk}_N(x)>0$, $y\notin\{x,*_N\}$,
    $\operatorname{rk}_N(y)\le\operatorname{rk}_N(x)$, and $\xi_y^{(N)}=\xi_x^{(N)}$, every
    $p\in B_x^{(N)}$ and $q\in R_y^{(\infty,N)}$ satisfy
    \begin{align}
        c_1\tau_N(x)^{\rho_0} & \le\|p-q\|_\infty.
        \label{eq:area_law_birth_separation_common_stage}
    \end{align}
\end{theorem}

\begin{proof}
    \leanok
    \uses{def:area_law_repair_birth_rank,
        def:area_law_repair_characteristic_length,
        thm:area_law_repair_final_retention}
    By Definition~\ref{def:area_law_repair_birth_rank}, $\operatorname{rk}_j(i)\le j+1$ for all
    $i\in I_j$, and Theorem~\ref{thm:area_law_repair_final_retention} gives
    $R_i^{(\infty,j)}=R_{\iota_{j,j+1}(i)}^{(\infty,j+1)}$. Induct on $N$: at $N=0$,
    $\operatorname{rk}_0(x)=1$ makes $x$ an initial belt with $\tau_0(x)=t_k$, so the initial bound
    applies; at $N=j+1$, if $x$ is fresh then $\tau_{j+1}(x)$ is the removal half-side and the
    replacement bound applies, if both $x,y$ are old then retention reduces to generation $j$,
    and $x$ old with $y$ fresh is impossible since $\operatorname{rk}_{j+1}(x)\le j+1<j+2=\operatorname{rk}_{j+1}(y)$.
\end{proof}

When $K_N=\varnothing$, Theorem~\ref{thm:area_law_repair_final_current_empty} gives
$R_y^{(\infty,N)}=R_y^{(N)}$.
